\scanpage{179}
\mspage{92}
% L01
X \uline{Antipodismes généralisés}, et \uline{présentation autoduale}.

% L02–07
.1. Soit $S$ un espace homogène sous un groupe \unsure{$G$|$G_0$}. Disons que deux points de $S$ sont \inserted{$G$-}\uline{équivalents} s'ils ont même stabilisateur \unsure{sous|dans} $G$, \deleted{\unsure{donc}} i.e. s'ils sont dans une même fibre \deleted{pour} \inserted{de} l'application
\[
\begin{aligned}
S&\longrightarrow\operatorname{Ssgr}(G)\\
s&\longmapsto G_s=\{g\in G\mid g.s=s\}.
\end{aligned}\tag*{(1)}
\]

% L08–11
\deleted{Les} \inserted{l'ens des} classes de $G$-équivalence de points \deleted{de} de l'espace homogène $S$ s'identifie, donc, \deleted{aux \unsure{à} l'ensemble} \inserted{en tant qu'espace où $G$ opère,} à une cl. de \inserted{conjugaison} de sous-groupes de $G$ -- savoir l'ensemble des \deleted{\unsure{$s$}} $G_s$.

% L12–18
Soit $Z\subset\unsure{\mathfrak S_S|\mathscr S_S}$ le commutant de $G$, alors $Z$ opère \uline{librement} sur $S$, et les classes de $G$-équivalence sont les $Z$-orbites dans $S$. Si $s_0\in S$, il y a donc corr. 1-1 entre l'ensemble des $s\in S$ qui sont $S$-équivalents à $S$, \deleted{entre} \deleted{et à \uline{l'orbite}} i.e. l'orbite de $s_0$ sous $Z$, et l'ens des $f\in\deleted{\unsure{G}} Z$, via l'application
\[
\begin{aligned}
f&\longmapsto f.s_0\\
Z&\xrightarrow{\sim} Z.s_0=\text{ens. des pts de }S\;G\text{-équivalents à }s_0.
\end{aligned}\tag*{(2)}
\]

% L19–23
Notons qu'on a
\[
\begin{aligned}
Z.s_0\subset S^{(G_{s_0})}&\overset{\mathrm{déf}}{=}\text{ens. des points de }S\\
&\quad\deleted{\text{\unsure{invariants|fixés}}}\inserted{\text{fixés}}\ \text{par le stabilisateur de }S,\\
&\quad\text{i.e. satisfaisant }G_s\subset G_{s_0}
\end{aligned}\tag*{(3)}
\]
Notons que le choix de $s_0$ comme ``origine'' permet d'interpréter $S$

% archive sheet number: 178
